% Problem 1
\begin{problem}{Properties of absolute value and distance}
	Verify the basic properties of absolute value and distance. Let $x, y, z \in \mathbb{Q}.$
\end{problem}
\begin{multicols}{2}
\begin{proposition}
    Non-degeneracy of absolute value. We have $\lvert x \rvert \geq 0.$ 
    Also, $\lvert x \rvert = 0$ if and only if $x = 0.$
\end{proposition}
\begin{proof}
    d
\end{proof}

\begin{proposition}
    (Triangle inequality for absolute value) We have $\lvert x + y \rvert \leq \lvert x \rvert + \lvert y \rvert.$
\end{proposition}
\begin{proof}
    s
\end{proof}

\begin{proposition}
    We have the inequalities $-y \leq x \leq y$ if and only if $y \geq \lvert x \rvert x.$
    In particular, we have $-\lvert x \rvert x leq x \leq \lvert x \rvert$.
\end{proposition}
\begin{proof}
    s
\end{proof}

\begin{proposition}
    (Multiplicativity of absolute value) We have $\lvert xy \rvert = \lvert x \rvert \lvert y \rvert.$ 
    In particular, $\lvert -x \rvert = \lvert x \rvert.$
\end{proposition}
\begin{proof}
    s
\end{proof}

\begin{proposition}
    (Non-degeneracy of distance) We have $d(x, y) \geq 0.$
    Also, $d(x, y) = 0$ if and only if $x = y.$
\end{proposition}
\begin{proof}
    smth
\end{proof}

\begin{proposition}
    (Symmetry of distance) We have $d(x, y) = d(y, x).$
\end{proposition}
\begin{proof}
    sh 
\end{proof}

\begin{proposition}
    (Triangle inequality for distance) We have $d(x, z) \leq d(x, y) + d(y, z).$
\end{proposition}

\begin{proof}
    chu
\end{proof}
\end{multicols}

% Problem 2
\begin{problem}{Properties of $\epsilon$-closeness}
Let $x, y, z, w \in \mathbb{Q}$. Then, the following properties hold:
\end{problem}
\begin{multicols}{2}
\begin{proposition}
    The number $x$ is $\varepsilon$-close to $y$ for every $\varepsilon > 0$ if and only if  $x = y$.
\end{proposition} 
\begin{proof}
    s
\end{proof}

\begin{proposition}
    Let $\varepsilon > 0.$ 
    If $x$ is $\varepsilon$-close to $y$, then $y$ is $\varepsilon$-close to $x$.
\end{proposition}
\begin{proof}
    s
\end{proof}

\begin{proposition}
    Let $\varepsilon, \delta > 0.$ If $x$ is $\varepsilon$-close to $y$, and $y$ is $\delta$-close to $z$, then $x$ and $z$ are $(\varepsilon + \delta)$-close. 
\end{proposition}
\begin{proof}
    chuchu
\end{proof}

\begin{proposition}
    Let $\varepsilon > 0.$ If $x$ and $y$ are $\varepsilon$-close, they are also $\varepsilon'$-close for every $\varepsilon' > \varepsilon$.
\end{proposition}
\begin{proof}
    chuchu
\end{proof}

\begin{proposition}
    Let $\varepsilon > 0.$ If $y$ and $z$ are both $\varepsilon$-close to $x$ and $w$ is between $y$ and $z$ (i.e., $y \leq w \leq z$ or $z \leq w \leq y$), then $w$ is also $\varepsilon$-close to $x$.
\end{proposition}
\begin{proof}
    chuchu
\end{proof}

\begin{proposition}
    Let $\varepsilon > 0.$ If $x$ and $y$ are $\varepsilon$-close, and $z$ is non-zero, then $xz$ and $yz$ are $\varepsilon \lvert z \rvert$-close.
\end{proposition}
\begin{proof}
    
\end{proof}

\begin{proposition}
    Let $\varepsilon, \delta > 0.$ Suppose $x$ and $y$ are $\varepsilon$-close, and $z$ and $w$ are $\delta$-close. Then $xz$ and $yw$ are $(\varepsilon \lvert z \rvert + \delta \lvert x \rvert + \varepsilon \delta)$-close.
\end{proposition}
\end{multicols}

% Probelm 4
\begin{problem}{Properties of Exponentiation I}
Let $x, y$ be rational numbers. Let $n, m$ be natural numbers. The following properties hold.
\end{problem}

\begin{multicols}{2}
\begin{proposition} We have $x^n x^m = x^{n + m}$.
\end{proposition}
\begin{proof}
    chuchu
\end{proof}

\begin{proposition} We have $(x^n)^m=x^{nm}.$
\end{proposition}
\begin{proof}
    chuchu
\end{proof}

\begin{proposition} $(xy)^n=x^n y^n.$
\end{proposition}
\begin{proof}
    chuchu
\end{proof}
\begin{proposition}
    Suppose $n \geq 0.$ Then, we have $x^n = 0$ if and only if $x = 0.$
\end{proposition}
\begin{proof}
    suppose
\end{proof}

\begin{proposition}
    If $x \geq y \geq 0$, then $x^n \geq y ^ n \geq 0$. 
\end{proposition}
\begin{proof}
    suppose
\end{proof}

\begin{proposition}
    If $x > y \geq 0$ and $n > 0$, then $x ^n > y^n \geq 0.$
\end{proposition}
\begin{proof}
    suppose
\end{proof}
\begin{proposition}
    We have $\lvert x^n \rvert = \lvert x \rvert ^n$.
\end{proposition}
\begin{proof}
    suppose
\end{proof}
\end{multicols}


% Probelm 5
\begin{problem}{Properties of Exponentiation II}
Let $x, y$ be rational numbers. Let $n, m$ be integers. The following properties hold.
\end{problem}

\begin{multicols}{2}
\begin{proposition} We have $x^n x^m = x^{n + m}$.
\end{proposition}
\begin{proof}
    chuchu
\end{proof}

\begin{proposition} We have $(x^n)^m=x^{nm}.$
\end{proposition}
\begin{proof}
    chuchu
\end{proof}

\begin{proposition} $(xy)^n=x^n y^n.$
\end{proposition}
\begin{proof}
    chuchu
\end{proof}

\begin{proposition}
    If $x \geq y > 0$, then $x^n \geq y^n > 0$ if $n$ is positive OR $0 < x^n \leq y^n$ if $n$ is negative.
\end{proposition}
\begin{proof}
    suppose
\end{proof}

\begin{proposition}
    If $x, y > 0,$ $n \neq 0,$ and $x^n = y^n$, then $x = y$.
\end{proposition}
\begin{proof}
    suppose
\end{proof}

\begin{proposition}
    We have $\lvert x^n \rvert = \lvert x \rvert ^n$.
\end{proposition}
\begin{proof}
    suppose
\end{proof}
\end{multicols}

\begin{problem}{Exponentiation grows faster than linear}
    Prove that $2^N \geq N$ for all positive integers $N$.
\end{problem}
\begin{proof}
    chuchu
\end{proof}